% Fixed preamble for transcribed notes. Compiled with XeLaTeX.
% The model only ever writes the document body; everything it may use is defined here.
\documentclass[11pt,a4paper]{article}

% --- Language and fonts --------------------------------------------------------
\usepackage{fontspec}               % UTF-8 input: å, ä, ö work directly
\usepackage[swedish,english]{babel} % last option (english) is the default
\usepackage{csquotes}

% --- Page layout ----------------------------------------------------------------
\usepackage[margin=2.2cm,marginparwidth=2cm]{geometry}
\usepackage{parskip}
\usepackage{enumitem}
\usepackage{multicol}               % formula sheets

% --- Maths ----------------------------------------------------------------------
\usepackage{amsmath,amssymb,mathtools,amsthm}  % amsthm: proof environment
\usepackage{physics}                % \vb, \dv, \pdv, \abs, \norm, \ket, \bra
\usepackage{bm}
\usepackage{cancel}                 % \cancel{x}: struck-out factors

% --- Science --------------------------------------------------------------------
\usepackage{siunitx}                % \SI, \si, \num, \qty
\sisetup{per-mode=symbol}
\usepackage[version=4]{mhchem}      % \ce{...}

% --- Graphics, colour, boxes ------------------------------------------------------
\usepackage{graphicx}
\usepackage{tikz}                   % redrawn diagrams (drawings/<id>.tex)
\usetikzlibrary{arrows.meta,decorations.markings,patterns}
\tikzset{>=Stealth,
  contour/.style={thick,blue!60!black},
  arrow at/.style={postaction={decorate},decoration={markings,mark=at position #1 with {\arrow{Stealth}}}},
  axis/.style={->,thin},
  pt/.style={circle,fill,inner sep=1.1pt},
  sing/.style={font=\small}}
\usepackage{xcolor}
\usepackage[most]{tcolorbox}
\usepackage{hyperref}
\hypersetup{colorlinks=true,linkcolor=blue!50!black}

% Box titles: "Definition" or "Definition: Title" when a title is given.
\newcommand{\nbtitle}[2]{#1\if\relax\detokenize{#2}\relax\else: #2\fi}

% Box names and the decimal marker follow the page language.
\newcommand{\definitionname}{Definition}
\newcommand{\theoremname}{Theorem}
\newcommand{\examplename}{Example}
\newcommand{\notename}{Note}
\addto\captionsenglish{%
  \renewcommand{\definitionname}{Definition}\renewcommand{\theoremname}{Theorem}%
  \renewcommand{\examplename}{Example}\renewcommand{\notename}{Note}%
  \sisetup{output-decimal-marker={.}}}
\addto\captionsswedish{%
  \renewcommand{\definitionname}{Definition}\renewcommand{\theoremname}{Sats}%
  \renewcommand{\examplename}{Exempel}\renewcommand{\notename}{Anmärkning}%
  \sisetup{output-decimal-marker={,}}}  % 9,82 stays 9,82 on Swedish pages

\newtcolorbox{definition}[1][]{enhanced,breakable,colback=blue!4,colframe=blue!55!black,
  fonttitle=\bfseries,title={\nbtitle{\definitionname}{#1}}}
\newtcolorbox{theorem}[1][]{enhanced,breakable,colback=red!4,colframe=red!55!black,
  fonttitle=\bfseries,title={\nbtitle{\theoremname}{#1}}}
\newtcolorbox{example}[1][]{enhanced,breakable,colback=green!4,colframe=green!45!black,
  fonttitle=\bfseries,title={\nbtitle{\examplename}{#1}}}
\newtcolorbox{note}[1][]{enhanced,breakable,colback=gray!6,colframe=gray!60!black,
  fonttitle=\bfseries,title={\nbtitle{\notename}{#1}}}

% List of theorems: the build step adds \theoremindex{title} to each theorem box.
\makeatletter
\newcommand{\theoremindex}[1]{\phantomsection\addcontentsline{thm}{subsection}{#1}}
\newcommand{\listoftheorems}{\section*{Satser och propositioner}\@starttoc{thm}}
\makeatother

% A framed section with only the page's own title (or none), for frames without a label.
\newtcolorbox{framed}[1][]{enhanced,breakable,colback=white,colframe=red!50!black,
  title={#1},fonttitle=\bfseries,coltitle=red!50!black,attach title to upper={\ }}

% --- Transcription helpers --------------------------------------------------------
% \unsure{...}: text the transcriber could not read with confidence (highlighted).
\newcommand{\unsure}[1]{\relax\ifmmode\colorbox{yellow!55}{$\displaystyle #1$}\else\colorbox{yellow!55}{#1}\fi}

% \rattat{new}{old}: a slip in the original notes, corrected at the student's request. Prints the
% corrected text with a green dagger; the original wording stays in the source.
\newcommand{\rattat}[2]{\relax\ifmmode{#1}{}^{\textcolor{green!50!black}{\dagger}}\else#1\textsuperscript{\textcolor{green!50!black}{\dag}}\fi}

% \notefigure[width]{id}{caption}: a diagram cropped from the original page. The build step
% fills in the width (share of \linewidth) from the drawing's width on the original page.
% A redrawn version in drawings/<id>.tex (TikZ) is used instead of the crop when present.
\newcommand{\notefigure}[3][0.5]{%
  \begin{center}
    \IfFileExists{drawings/#2.tex}%
      {\input{drawings/#2.tex}}%
      {\IfFileExists{figures/#2.png}%
        {\includegraphics[width=#1\linewidth,height=0.35\textheight,keepaspectratio]{figures/#2.png}}%
        {\fbox{\parbox{0.5\linewidth}{\centering\small [Figure #2 missing]}}}}%
    \if\relax\detokenize{#3}\relax\else\par\small\emph{#3}\fi
  \end{center}}

% \pagebreakmarker: separates pages of the original notes.
\newcommand{\pagebreakmarker}[1]{\par\medskip\noindent{\color{gray}\rule{\linewidth}{0.3pt}\\[-0.4em]\hfill\scriptsize\MakeLowercase{\pagename} #1}\par\medskip}
